\documentclass[10pt,a4paper]{amsart}
\usepackage[margin=19mm]{geometry}
\usepackage{amsmath,amssymb,mathtools}
\usepackage[scr=rsfs,cal=euler]{mathalfa}
\usepackage{xfrac}
\usepackage[unicode,hidelinks,bookmarks=false]{hyperref}
\newcommand{\ie}{i.e.\ }
\newcommand{\cf}{cf.\ }
\newcommand{\resp}{respectively\ }
\newcommand{\Subsection}{\subsection}
\providecommand{\nobreakdash}{\nobreak\hbox{-}}
\setlength{\emergencystretch}{2em}
\multlinegap0pt
\usepackage{bm}
\usepackage{bbm}
\usepackage{dsfont}
\usepackage{xspace}
\usepackage{cancel}
\usepackage{mathtools}
\usepackage{amsthm}
\makeatletter
\@ifundefined{theorem}{\newtheorem{theorem}{Theorem}}{}
\@ifundefined{assumption}{\newtheorem{assumption}[theorem]{Assumption}}{}
\@ifundefined{lemma}{\newtheorem{lemma}[theorem]{Lemma}}{}
\@ifundefined{proposition}{\newtheorem{proposition}[theorem]{Proposition}}{}
\makeatother
\DeclareMathOperator{\diver}{div}
\DeclarePairedDelimiter{\norm}{\lVert}{\rVert}
\newcommand{\Ac}{\mathcal{A}}
\newcommand{\Af}{{A}}
\newcommand{\Cregz}{C_{z,\text{reg}}}
\newcommand{\Hc}{\mathcal{H}}
\newcommand{\om}{\Omega}
\title[Reliable eigenspace error estimation using source error esti]{Reliable eigenspace error estimation using source error estimators}
\author{Jay Gopalakrishnan, Gabriel Pinochet-Soto}
\date{Source: arXiv:2601.08051v2 (CC BY 4.0)}
\begin{document}
\maketitle

\noindent\emph{Source and licence.} Reliable eigenspace error estimation using source error estimators. Version arXiv:2601.08051v2 (CC BY 4.0), distributed under the Creative Commons Attribution 4.0 International licence, as recorded in the arXiv licence metadata for this exact version and shown on the abstract page https://arxiv.org/abs/2601.08051v2. Full rights notice: licenses/arxiv-2601.08051-CC-BY-4.0.txt.\par\smallskip

\noindent\textit{Editorial scope.} Counted unit: the Proposition whose source label is \texttt{None} in the article (source0.tex lines 1921--1924, source label \texttt{None}), delivered together with the article's own complete original proof of it (source0.tex lines 1925--1980, one \texttt{proof} environment of 56 lines). No mathematics is written, repaired or re-derived: statement and proof are the article's own text line for line. Required context retained uncounted: asm:resolvent-approx (lines 1318--1329); eq:regularity (lines 1704--1708); lem:Laplace-elliptic' (lines 1848--1861). Every definition, display and label of those lines is retained unchanged. Omitted from this package and not redistributed: 1--1317, 1330--1703, 1709--1847, 1862--1920, 1981--3021. Nothing inside a retained excerpt is omitted or altered: every retained body is delivered exactly as the article's lines are written, so no omission receipt of any kind accompanies this package. Citations: the retained excerpts carry no citation command at all, so no bibliography entry of the article is transcribed and no reference is redistributed. Reference adaptation: none was needed and none was made; every reference inside the delivered unit resolves inside it, so no unresolved reference or citation is produced. Printed slips kept literal: no printed word, symbol, hypothesis or number of the counted statement or of its proof was altered. Layout and class: the article's own class is \texttt{amsart} with packages including fontenc, geometry, todonotes, amsmath, amssymb, amsthm, amsfonts, amsaddr, mathtools, mathrsfs, graphicx, hyperref, xcolor, tikz; the wrapper is the lane's pinned \texttt{amsart} editorial wrapper at 10pt on A4, which restates the theorem-like environments the retained text needs and adds the article's own macro definitions that the retained text needs (\texttt{\textbackslash Ac}, \texttt{\textbackslash Af}, \texttt{\textbackslash Cregz}, \texttt{\textbackslash Hc}, \texttt{\textbackslash diver}, \texttt{\textbackslash norm}, \texttt{\textbackslash om}), with extra wrapper packages (bm, bbm, dsfont, xspace, cancel, mathtools). The delivered numbering is the unit's own. Assets: no retained span contains a figure environment, a graphics-inclusion command, a PDF-page inclusion, a table or a TikZ picture; no image file is copied into this package, the package declares no asset dependency and no image file is redistributed with it. Licence: version arXiv:2601.08051v2 of this article is distributed under the Creative Commons Attribution 4.0 International licence, as recorded in the arXiv licence metadata for this exact version and shown on the abstract page https://arxiv.org/abs/2601.08051v2. A full rights notice is delivered at licenses/arxiv-2601.08051-CC-BY-4.0.txt. Characters: every retained excerpt is pure ASCII, so the delivered engine, XeLaTeX with its default Latin Modern text font, reports no missing character.
\begin{assumption}[On Resolvent Approximations]
  \label{asm:resolvent-approx}
  We assume that there are
  finite-dimensional subspaces $V_h \subset V$ (indexed by a
  discretization parameter $h$ approaching zero) and finite-rank
  operators $R_h(z) : \Hc \to V_h$ such that
  \begin{equation}
    \label{eq:1}
    \lim_{h \to 0} \| R_h(z_k) - R(z_k) \|_V = 0
  \end{equation}
  for every $k = 1, \dots, N.$
\end{assumption}

\begin{equation}
  \label{eq:regularity}
  \| R(z) f \|_{H^{1+s}(\om)} \le \Cregz \| f \|_{\Hc},
  \qquad \text{ for all }  f \in \Hc \text { and } z \in \rho(\Ac).
\end{equation}

\begin{lemma}
  \label{lem:Laplace-elliptic'}
  For any \(z \in \rho(\Ac)\), there exists \(\gamma_z > 0\) such that, for all
  \(u \in \mathring{H}^1(\Omega)\) and \(q \in H(\diver)\), we have
  \begin{equation}
    \label{eq:Laplace-elliptic'}
    \norm{q}_{L_2% (\Omega)^2
    }^2 +
    \norm{u}_{H^1(\Omega)
    }^2 
    \leq \gamma_z \norm{\Af_z (q, u)}_{L_2}^2.
  \end{equation}
  Moreover, the operator \(\Af_z\) is a continuous bijection.
\end{lemma}

\begin{proposition}
  The FOSLS resolvent approximation $R_h(z)$ converges in operator norm on $V$,
  i.e., Assumption~\ref{asm:resolvent-approx} holds.
\end{proposition}

\begin{proof}
  Let $f \in V$,  $u = R(z) f$,  and $u_h = R_h(z) f$.
  Since $V_h \subset V$, by Galerkin orthogonality
  \begin{align*}
    \norm*{\Af_z (q - q_h, u - u_h)}_{L_2}^2
    & =
      \left(
      \Af_z(q - q_h, u - u_h), \Af_z(q - r_h, u - w_h)
      \right)_{L_2}
    \\
    & \lesssim
      \| \Af_z(q - q_h, u - u_h)\|_{L_2}
      \Big(
      \| q - r_h \|_{H(\diver)}  + \| u - w_h \|_{H^1(\om)} 
      \Big).
  \end{align*}
  By standard finite element best approximation estimates for the Lagrange
  and Raviart-Thomas spaces, this leads to 
  \begin{equation}    
    \label{eq:7-Az}
    \norm*{\Af_z (q - q_h, u - u_h)}_{L_2}
    \lesssim h^r |u |_{H^{r+1}(\om} + h^r |q|_{H^r(\om)} + h^r | \diver q |_{H^r(\om)}
  \end{equation}
  for some $0 < r \le 1$. Choosing $r\le s$, where $s$ is as
  in~\eqref{eq:regularity}, the seminorms on the right hand side above can
  be bounded as follows:
  \begin{gather*}
    |u|_{H^{r+1}(\om)} + |q|_{H^{r}(\om)}
    \lesssim \| u \|_{H^{s+1}(\om)} \lesssim
      \| f \|_{\Hc} \lesssim \| f \|_V,
    \\
    |\diver q|_{H^{r}(\om)}
    = |z u  - f|_{H^r(\om)} \lesssim \| f \|_{H^r(\om)} \lesssim \| f \|_V,
  \end{gather*}
  where in the last step we have used that $r \le 1$ as well as the
  first bound for~$u$.  Using these estimates in~\eqref{eq:7-Az},
  \[
    \norm*{\Af_z (q - q_h, u - u_h)}_{L_2}
    \lesssim\, h^r \| f \|_V.
  \]
  By Lemma~\ref{lem:Laplace-elliptic'},
  $    \norm*{u - u_h}_{H^1(\Omega)}^2
  + \norm*{q - q_h}_{L_2}^2
  \leq
  \gamma_z \norm*{\Af_z (q - q_h, u - u_h)}_{L_2}^2,
  $
  so we conclude that for any $f \in V$, 
  \begin{align}
    \label{eq:8}
    \| R(z) f - R_h(z) f\|_V
    & = \| u - u_h \|_{H^1(\om)}
      \le \gamma_z^{1/2} \| \Af_z( q - q_h, u - u_h) \|_{L_2} 
      \lesssim h^r \| f \|_V,    
  \end{align}
  thus verifying Assumption~\ref{asm:resolvent-approx} as $h \to 0$. 
\end{proof}

\end{document}
